\documentclass[12pt]{article}
\usepackage[utf8]{inputenc}
\usepackage[T1]{fontenc}
\usepackage[margin=1in]{geometry}
\usepackage{setspace}
\usepackage{xcolor}
\usepackage[colorlinks=true, urlcolor=blue!60!black, linkcolor=blue!60!black]{hyperref}
\setlength{\parindent}{2em}
\setlength{\parskip}{1.25em}
\renewcommand{\baselinestretch}{1.0}
\begin{document}
\begin{center}
{\LARGE \textit{303 Creative LLC v.\ Elenis}, 600 U.S.\ 570 (2023)\par}
\vspace{0.8em}
{\large Case Brief --- Four-Page Edition\par}
\end{center}
\vspace{1em}

\section{Memory Jogger}

The Free Speech Clause bars Colorado from using its public accommodations law to compel a commercial website designer to create customized, expressive wedding websites conveying messages that conflict with her beliefs (\href{https://supreme.justia.com/cases/federal/us/600/21-476/}{Justia}).

\section{Detailed Case Facts}

Lorie Smith, sole member-owner of 303 Creative LLC, a Colorado graphic and website design business, wanted to expand into wedding websites telling each couple's love story, but only for opposite-sex weddings, because of her Christian belief that marriage is between one man and one woman. She also wanted to post a statement explaining that refusal. The Colorado Anti-Discrimination Act (CADA) bars a public accommodation from denying ``the full and equal enjoyment'' of its services based on sexual orientation (the Accommodation Clause) and from publishing that a person's patronage is unwelcome (the Communication Clause). The parties stipulated that Smith will serve LGBTQ+ clients in non-conflicting projects, that every website she designs is an ``original, customized'' and ``expressive'' creation, and that many other companies offer custom website design. She had not yet designed any wedding website when she sued. The suit followed \textit{Masterpiece Cakeshop} (2018), decided narrowly on Free Exercise grounds.

\section{Procedural History}

Smith brought a pre-enforcement challenge. The district court granted summary judgment to Colorado. A divided Tenth Circuit found standing and agreed the websites were ``pure speech'' triggering strict scrutiny, but held 2--1 that CADA survived because Smith's services were ``unique''; Chief Judge Tymkovich dissented. The Supreme Court granted certiorari limited to the Free Speech question. On remand, the district court entered a permanent injunction.

\section{Judicial Votes}

Reversed, \textbf{6--3} (decided June 30, 2023).
\begin{itemize}
\item \textbf{Majority:} Justice Gorsuch, joined by Chief Justice Roberts and Justices Thomas, Alito, Kavanaugh, and Barrett.
\item \textbf{Concurrences:} None.
\item \textbf{Dissent:} Justice Sotomayor, joined by Justices Kagan and Jackson.
\end{itemize}

\section{Holding}

A state may not apply its public accommodations law to compel a business that creates customized, expressive ``pure speech'' to design works conveying messages the creator disagrees with, even when it sells those services for profit to the public through a corporate entity. The holding rests on the stipulations; the Court declined to announce a general test for which other commercial services qualify as ``expressive activity.''

\section{Analysis of Opinions}

\textbf{Majority (Gorsuch).} Relying on \textit{Barnette}, \textit{Hurley}, and \textit{Dale}, the Court held that government may neither force a person to speak a state-preferred message nor use a public accommodations law to alter a speaker's expressive content. Smith's websites are pure speech and are \emph{her} speech, because she vets projects and authors original words and artwork; combining her voice with the couple's does not forfeit protection. States have a compelling interest in eliminating discrimination in public accommodations, but those laws cannot override the Constitution, and treating an artist's unique talent as a reason for coercion would invert the First Amendment. Payment and LLC status do not strip protection; Smith's objection is to a message, not a status; and \textit{FAIR} involved only an incidental burden. The Court warned that Colorado's rule would let states force a gay website designer to build a site opposing same-sex marriage.

\textbf{Dissent (Sotomayor).} Public accommodations laws, from the common-law duty of businesses that ``hold themselves out'' to the Civil Rights Act of 1964, secure equal access and equal dignity. The Court had consistently rejected First Amendment exemptions for commercial refusals of service (\textit{Heart of Atlanta Motel}, \textit{Katzenbach v.\ McClung}, \textit{Runyon}, \textit{Hishon}, \textit{Roberts}). CADA regulates the conduct of refusing service based on status, and any burden on speech is ``incidental'' under \textit{O'Brien} and \textit{FAIR}. Recasting the service as ``opposite-sex wedding websites'' collapses status into message, and the logic reaches interracial couples and disabled parents.

\section{Examples: Future Applications}

\begin{enumerate}
\item \textbf{Same-Side: Bespoke Wedding Photographer.} A photographer who curates custom visual narratives refuses a same-sex wedding but takes portraits for LGBTQ+ clients. Compelling the shoot compels pure speech; the Second Circuit found such a claim plausible in \textit{Emilee Carpenter, LLC v.\ James} (2024).
\item \textbf{Same-Side: Freelance Speechwriter.} An atheist speechwriter open to the public refuses to write a speech celebrating a religious revival. Original prose is quintessential pure speech.
\item \textbf{Opposite-Side: Reception Venue.} A ballroom that rents space and serves banquet food refuses a same-sex reception. Space and food are ordinary commercial goods, so CADA validly regulates the refusal.
\item \textbf{Opposite-Side: Template Print Shop.} A shop that prints customers' own PDFs or swaps names into template cards sells a standardized product, so its refusal is regulable status-based conduct.
\item \textbf{Fence-Sitter: Custom Cake Baker or Florist.} Cakes and floral arrangements do not automatically qualify as pure speech, and the Court left open when partly functional crafts become ``expressive activity'' attributed to the creator.
\end{enumerate}

\section{Critique}

\textbf{Doctrinal ambiguity.} Scholars note that the majority never specifies a tier of scrutiny, applying a categorical bar against compelled pure speech. Robert Post criticizes its ``impossible abstraction,'' and commentators fault the lack of a workable attribution test. \textbf{Contrived record.} Because Smith had never designed a wedding website, broad stipulations let the two sides ``talk past each other'' (David Cole), and the ``Stewart'' inquiry controversy fed charges of a contrived pre-enforcement record.

\textbf{Egalitarian and originalist views.} Progressive scholars argue the Court carved a deregulatory hole in civil rights law; Michael Dorf notes that owners, not employees, get the shield; David Cole reads the case narrowly. Robert G.\ Natelson argues an originalist analysis of the press freedom would have made the case easier for Smith.

\section{Key Quotations}

\begin{quote}
``When a state public accommodations law and the Constitution collide, there can be no question which must prevail.''
\end{quote}
\noindent --- Justice Gorsuch, for the Court.

\begin{quote}
``Were the rule otherwise, the better the artist, the finer the writer, the more unique his talent, the more easily his voice could be conscripted to disseminate the government's preferred messages. That would not respect the First Amendment; more nearly, it would spell its demise.''
\end{quote}
\noindent --- Justice Gorsuch, for the Court.

\begin{quote}
``Today, the Court, for the first time in its history, grants a business open to the public a constitutional right to refuse to serve members of a protected class.''
\end{quote}
\noindent --- Justice Sotomayor, dissenting.

\end{document}
